% Previous standalone explanations, retained for reference; not included in the presentation.
\begin{frame}[shrink=8]{Track partitions: configurable names}
        \small Groups: \texttt{TrackSelection1},
        \texttt{TrackSelection2}, \ldots\par
        \medskip
        \footnotesize
        \renewcommand{\arraystretch}{1.0}
        \begin{tabular}{p{.53\textwidth}p{.39\textwidth}}
                \toprule
                Configurable & Meaning\\
                \midrule
                \texttt{pdgCodeAbs} & PDG code for the mass hypothesis.\\
                \texttt{chargeSign}, \texttt{chargeAbs} & Sign and
                absolute charge.\\
                \texttt{ptMin}, \texttt{ptMax} & Transverse-momentum range.\\
                \texttt{etaMin}, \texttt{etaMax} & Pseudorapidity range.\\
                \texttt{phiMin}, \texttt{phiMax} & Azimuthal-angle range.\\
                \texttt{maskLowMomentum}, \texttt{maskHighMomentum} &
                Required bits below/above threshold.\\
                \texttt{pidThres} & Momentum threshold (GeV/$c$).\\
                \texttt{rejectionMaskLowMomentum},\newline\texttt{rejectionMaskHighMomentum}
                & Veto bits below/above threshold.\\
                \texttt{massMin}, \texttt{massMax} & TOF mass
                (optional process).\\
                \texttt{puritySlotMask} & Purity slots (FitVars only).\\
                \bottomrule
        \end{tabular}
        \medskip
        \footnotesize\textbf{Helium-3:} stored $p_{\mathrm T}/|Z|$ is
        scaled by \texttt{chargeAbs = 2}; most particles have $|Z|=1$.\par
        \textbf{\texttt{chargeSign = 0}:} keep both charge signs,
        useful in derived-to-derived production.
        \note{Names are copied from ConfTrackSelection in the local
                installation. Use the group prefix in full option names, for
                example TrackSelection1.ptMin. pdgCodeAbs sets a mass
                hypothesis rather than imposing PID. Stored signedPt is a
                signed rigidity-like quantity; the partition uses
                abs(chargeAbs * signedPt) for physical transverse momentum.
                For helium-3 and anti-helium-3 set chargeAbs=2. chargeSign=0
                removes the charge-sign restriction, useful when retaining
                both charges in derived-to-derived production. The
                FitVars-specific puritySlotMask controls QA output, not a
                partition cut. Partitions may overlap. Values must remain
        within what the producer retained.}
        \src{O2Physics/PWGCF/Femto/Core/trackBuilder.h;
        O2Physics/PWGCF/Femto/Core/partitions.h}
\end{frame}

\begin{frame}[shrink=8]{Example: $\Lambda$ partition configurables}
        \small Groups: \texttt{LambdaSelection1},
        \texttt{LambdaSelection2}, \ldots\par
        Example full name: \texttt{LambdaSelection1.massMin}.
        \medskip
        \footnotesize
        \renewcommand{\arraystretch}{1.3}
        \begin{tabular}{p{.43\textwidth}p{.49\textwidth}}
                \toprule
                Configurable & Meaning\\
                \midrule
                \texttt{pdgCodeAbs} & Mass hypothesis ($3122$ for $\Lambda$).\\
                \texttt{sign} & $+1$: $\Lambda$; $-1$:
                $\overline\Lambda$; $0$: both.\\
                \texttt{ptMin}, \texttt{ptMax} & Candidate
                transverse-momentum range.\\
                \texttt{etaMin}, \texttt{etaMax} & Candidate
                pseudorapidity range.\\
                \texttt{phiMin}, \texttt{phiMax} & Candidate
                azimuthal-angle range.\\
                \texttt{massMin}, \texttt{massMax} & Invariant-mass
                window or sideband.\\
                \texttt{mask} & Required topology and daughter-selection bits.\\
                \bottomrule
        \end{tabular}
        \medskip
        \small
        \textbf{Example: $\Lambda\to p\pi^-$.} Set \texttt{sign = 1};
        obtain \texttt{mask} from \texttt{hLambdaSelection} with CutCulator.
        \par\medskip
        \textbf{\texttt{sign = 0}:} keep both $\Lambda$ and
        $\overline\Lambda$, useful in derived-to-derived production.
        \note{Names are copied from ConfLambdaSelection and
                V0_DEFAULT_SELECTIONS. Lambda is an example of the partition
                concept. The consumer uses one V0 mask rather than track
                low/high momentum masks. sign distinguishes Lambda from
                anti-Lambda, not electric charge. Topology bits include
                pointing angle, daughter DCA, decay vertex and transverse
                radius. The mass range is limited by producer acceptance. For
        anti-Lambda metadata use hAntiLambdaSelection.}
        \src{O2Physics/PWGCF/Femto/Core/v0Builder.h;
        O2Physics/PWGCF/Femto/Core/partitions.h}
\end{frame}

\begin{frame}{From selected particles to a correlation}
        \begin{columns}[T]
                \begin{column}{.48\textwidth}
                        \begin{block}{Same event: $A(k^*)$}
                                Partition 1, event $i$\par
                                $\qquad\times$ partition 2, event
                                $i$\par\medskip
                                Retains the physical correlations.
                        \end{block}
                \end{column}
                \begin{column}{.48\textwidth}
                        \begin{block}{Mixed events: $B(k^*)$}
                                Partition 1, event $i$\par
                                $\qquad\times$ partition 2, event
                                $j\ne i$\par\medskip
                                Build a reference using similar events.
                        \end{block}
                \end{column}
        \end{columns}
        \begin{itemize}
                \item Mix within suitable vertex and
                        multiplicity/centrality classes.
                \item Apply pair cleaning and close-pair rejection as
                        appropriate.
                \item Exclude self-pairs and handle duplicate
                        combinations for identical samples.
                \item Fill same- and mixed-event pair distributions.
        \end{itemize}
        \note{Mixed-event pairs provide a reference, not an automatic
                correction of every detector or event-class effect. Inspect
                mixing QA. Pair cleaning matters for shared daughters in
                composite-particle channels. Purity/fraction corrections and
        physics fits are later analysis steps.}
        \src{O2Physics/PWGCF/Femto/Tasks/femtoPairTrackTrack.cxx;
        O2Physics/PWGCF/Femto/Core/pairBuilder.h}
\end{frame}

\begin{frame}{One workflow, demonstrated locally and on Hyperloop}
        \centering
        \begin{tikzpicture}[node
                        distance=6mm,box/.style={draw=teal,rounded
                                corners,fill=pale,text width=2.7cm,minimum
                height=18mm,align=center},>=Stealth]
                \node[box] (a) {Produce\\\small AO2D $\to$ derived data};
                \node[box,right=of a] (b) {Choose\\\small QA file
                $\to$ CutCulator};
                \node[box,right=of b] (c) {Analyze\\\small partitions
                $\to$ pairs};
                \node[box,right=of c] (d) {Inspect\\\small QA and correlations};
                \draw[->,thick] (a)--(b);\draw[->,thick]
                (b)--(c);\draw[->,thick] (c)--(d);
        \end{tikzpicture}
        \medskip
        \begin{columns}[T]
                \begin{column}{.48\textwidth}
                        \begin{block}{Local demonstration}
                                A small input sample, explicit
                                task/configuration files, immediate
                                output inspection.
                        \end{block}
                \end{column}
                \begin{column}{.48\textwidth}
                        \begin{block}{Hyperloop demonstration}
                                Corresponding workflows and options
                                at production scale; inspect the
                                retrieved ROOT output with the same tools.
                        \end{block}
                \end{column}
        \end{columns}
        \medskip
        \raggedright\textbf{Keep together:} software revision,
        producer configuration, producer QA metadata, derived-data
        provenance and consumer settings.
        \note{Transition into the live demos; no exercises. This
                introductory deck does not claim that a Hyperloop train is
                already configured. 
                Local modifications and Hyperloop software versions must
        agree where schema, masks and task options matter.}
\end{frame}

% Archived standalone producer explanations; not included by main.tex.
\section{Configuring the producer}

\begin{frame}[shrink=8]{Producer configurables: three questions per object}
 \small
 \renewcommand{\arraystretch}{1.2}
 \begin{tabular}{p{.24\textwidth}p{.29\textwidth}p{.37\textwidth}}
 \toprule
 Group & Question & Track example\\
 \midrule
 \texttt{*Filters} & Which candidates enter the selection? & \texttt{TrackFilters.ptMin}\\
 \texttt{*Bits} & Which decisions are evaluated and encoded? & \texttt{TrackBits.tpcClustersMin}\\
 \texttt{*Tables} & Which output tables are filled? & \texttt{TrackTables.produceTracks}\\
 \bottomrule
 \end{tabular}
 \medskip
 \begin{block}{The same pattern across particle builders}
 Collisions, tracks, V0s, cascades and other candidates have their own configuration groups. The producer brings these builders together.
 \end{block}
 \textbf{Process switches are another layer:} activate the producer process matching the system, data/MC and input candidate types.
 \note{The wildcard names express the design pattern, not literal JSON keys. Check each builder for exact prefixes: the local collision prefilter prefix is CollisionFilter (singular), while the track prefix is TrackFilters. Enabling an output table alone does not activate the process that fills it.}
 \src{O2Physics/PWGCF/Femto/TableProducer/femtoProducer.cxx; O2Physics/PWGCF/Femto/Core/trackBuilder.h; O2Physics/PWGCF/Femto/Core/collisionBuilder.h}
\end{frame}

\begin{frame}[fragile,shrink=8]{Prefilters define the production acceptance}
 \small
 \begin{lstlisting}
 TrackFilters.ptMin  = 0.2
 TrackFilters.ptMax  = 6.0
 TrackFilters.etaMin = -0.9
 TrackFilters.etaMax =  0.9
 \end{lstlisting}
 These are the local default values shown as option names and values.
 \begin{itemize}\setlength{\itemsep}{2pt}
 \item Fixed bounds run before bit selections; failed standalone candidates are removed.
 \item Consumer kinematic cuts can narrow this acceptance later.
 \end{itemize}
 \begin{block}{Choose an envelope covering the planned analyses}
 Producing only $p_{\mathrm T}>0.5$ GeV/$c$ cannot support a later standalone-track analysis down to $0.2$ GeV/$c$.
 \end{block}
 \small Prefilter QA: \texttt{hTrackFilters}. Collision bounds use \texttt{CollisionFilter.*}, e.g. \texttt{CollisionFilter.vtxZMin}.
 \note{For tracks these bounds are strict comparisons. The standalone fill loop checks filters, then bit selections, then acceptance, and finally fills the output. Composite-candidate daughter tracks may enter via their own builder path, so avoid claiming all rows of the shared track table must obey the standalone prefilter. TrackBits.passThrough bypasses the bit acceptance gate, not this prefilter.}
 \src{O2Physics/PWGCF/Femto/Core/trackBuilder.h}
\end{frame}

\begin{frame}[fragile,shrink=8]{Bit selections define the stored working points}
 \begin{lstlisting}
 TrackBits.tpcClustersMin = [90, 100, 110]
 TrackBits.tpcProton      = ["-3;3", "-2.5;2.5"]
 TrackBits.requirePidProton = true
 \end{lstlisting}
 \small Illustrative option values, not a complete producer configuration.
 \begin{itemize}
 \item A vector configures several thresholds or ranges for one observable.
 \item TPC cluster quality is minimal: 90 is the baseline; 100 and 110 provide tighter decisions.
 \item \texttt{requirePidProton=true} marks configured proton PID selections as optional in the producer acceptance logic.
 \item With \texttt{false}, configured proton PID selections are neutral instead.
 \end{itemize}
 \begin{block}{Producer settings $\to$ metadata $\to$ consumer masks}
 Inspect \texttt{hTrackSelection} with CutCulator. The loosest quality bit may be omitted; use the stored mapping to choose tighter quality and PID requirements.
 \end{block}
 \note{The values above illustrate C++ vector contents, not literal full JSON serialization. The addSelection calls mark track quality as minimal and skip the loosest bit. PID ranges carry bits; requirePidProton controls their optional role, despite the potentially misleading word require. Several configured optional observables/species still form an OR at producer acceptance. Configure all desired species and tighter working points before production.}
 \src{O2Physics/PWGCF/Femto/Core/trackBuilder.h}
\end{frame}

\begin{frame}[fragile]{Tables separate the compact core from extra information}
 \begin{lstlisting}[language=C++]
 // Standard track-pair consumer joins matching rows:
 using FemtoTracks = o2::soa::Join<o2::aod::FTracks,
                                 o2::aod::FTrackMasks>;
 \end{lstlisting}
 \small
 \renewcommand{\arraystretch}{1.25}
 \begin{tabular}{p{.27\textwidth}p{.35\textwidth}p{.28\textwidth}}
 \toprule
 Content & Table examples & Producer switches\\
 \midrule
 Collision context + mask & \texttt{FCols}, \texttt{FColMasks} & \texttt{CollisionTables.*}\\
 Track kinematics + mask & \texttt{FTracks}, \texttt{FTrackMasks} & \texttt{TrackTables.*}\\
 Optional observables & \texttt{FTrackDcas}, \texttt{FTrackExtras}, \texttt{FTrackMass} & \texttt{TrackTables.*}\\
 PID / purity information & \texttt{FProtonPids}, \texttt{FTpcPurity}, \ldots & \texttt{TrackTables.*}\\
 \bottomrule
 \end{tabular}
 \medskip
 Collision indices link tracks to their event. Extension tables supply extra columns for the corresponding base rows.
 \note{An soa::Join here combines row-aligned columns rather than doing an arbitrary key-based SQL join. The base track table contains collision association and kinematics; the mask lives in its own extension table. Additional data increase the output size but enable studies beyond stored Boolean cuts. Composite channels also need daughter/index dependencies. Lite base tables are alternatives to full base tables in this installation; do not enable both representations simultaneously.}
 \src{O2Physics/PWGCF/Femto/DataModel/FemtoTables.h; O2Physics/PWGCF/Femto/Core/trackBuilder.h; O2Physics/PWGCF/Femto/Core/collisionBuilder.h}
\end{frame}

\begin{frame}{Table switches: auto, forced off, forced on}
 \small All \texttt{produce*} options in these table groups use integer settings:
 \medskip
 \renewcommand{\arraystretch}{1.45}
 \begin{tabular}{clp{.64\textwidth}}
 \toprule
 Value & Mode & Effect\\
 \midrule
 \texttt{-1} & Auto & Enable when the table is required in the current workflow.\\
 \texttt{0} & Forced off & Do not fill the table, even if a consumer needs it.\\
 \texttt{1} & Forced on & Fill the table independently of consumer demand.\\
 \bottomrule
 \end{tabular}
 \medskip
 \begin{columns}[T]
 \begin{column}{.48\textwidth}
 \begin{block}{Full workflow: producer + consumer}
 Leave tables on \textbf{auto}. The consumer's table inputs identify what the workflow needs.
 \end{block}
 \end{column}
 \begin{column}{.48\textwidth}
 \begin{block}{Producing derived data for later}
 \textbf{Force on} the tables to preserve. Future consumers are absent from the current workflow.
 \end{block}
 \end{column}
 \end{columns}
 \note{utils::enableTable first handles 1 and 0 explicitly, otherwise calls isTableRequiredInWorkflow. Use the documented -1 value for auto. Auto inspects the current workflow, not future analysis requirements. Disabling a required table can leave missing or empty inputs; it is not a safe way to optimize a consumer workflow. Enabling filling and selecting tables for the output writer are separate decisions.}
 \src{O2Physics/PWGCF/Femto/Core/femtoUtils.h}
\end{frame}

\begin{frame}[fragile,shrink=8]{Derived-data production: declare what to preserve}
 \begin{lstlisting}
 # Standard collision/track core for later track-pair analysis:
 CollisionTables.produceCollisions     = 1
 CollisionTables.produceCollisionMasks = 1
 TrackTables.produceTracks             = 1
 TrackTables.produceTrackMasks         = 1
 # Example extension for a later DCA study:
 TrackTables.produceTrackDcas          = 1
 \end{lstlisting}
 \small Option/value sketch; process selection and output-writer settings are additional configuration.
 \begin{enumerate}
 \item Enable the required producer process and any input dependencies.
 \item Force on the base, mask, extension and index-dependency tables needed later.
 \item Configure the derived-data writer to persist those tables.
 \item Keep producer QA for CutCulator; check the written table inventory.
 \end{enumerate}
 \textbf{Local and Hyperloop:} use the same table-selection logic; inspect the enabled-table logs and saved output in both cases.
 \note{This is the standard full-table track path, not a universal minimal dataset for every analysis. A purity study or composite channel needs more tables. Force on does not itself write an AO2D file; the writer must select/persist the produced tables. AnalysisResults.root holds QA histograms, while derived tables go to the configured AOD output. Full versus Lite representations are mutually exclusive in the local builders; Lite data need the corresponding conversion path before a standard full-table consumer. Avoid presenting this sketch as runnable shell syntax or a complete JSON configuration.}
 \src{O2Physics/PWGCF/Femto/Core/trackBuilder.h; O2Physics/PWGCF/Femto/Core/collisionBuilder.h}
\end{frame}


\begin{frame}{Code and tools used in this introduction}
        \small Source paths relative to \texttt{O2Physics/PWGCF/Femto/}:
        \begin{itemize}
                \item \texttt{TableProducer/femtoProducer.cxx}:
                        production entry point.
                \item \texttt{Core/trackBuilder.h}: track
                        configuration and selections.
                \item \texttt{Core/baseSelection.h},
                        \texttt{Core/selectionContainer.h}: roles,
                        acceptance and bit assignment.
                \item \texttt{Macros/cutculator.py},
                        \texttt{Macros/cutculator\_gui.py}: terminal
                        and GUI tools.
                \item \texttt{Core/partitions.h}: selection expressions.
                \item \texttt{Tasks/femtoPairTrackTrack.cxx}: consumer example.
                \item \texttt{Core/collisionBuilder.h},
                        \texttt{Core/femtoUtils.h}: table switches.
        \end{itemize}
        \medskip
        Source checkout: \texttt{\$HOME/alice/O2Physics}\par
        \medskip
        \href{https://github.com/AliceO2Group/O2Physics/tree/master/PWGCF/Femto}{Upstream:
        AliceO2Group/O2Physics, PWGCF/Femto}
        \note{The installed checkout is the primary reference. In
                particular, GUI export and TUI decode features described here
        are verified in the local Python source; upstream master may differ.}
\end{frame}
